\section*{Mock 4 · 10 points}
Model answers below each have three sentences. Accept other valid answers
that respect the stated assumptions.
\paragraph{1. The inflation-proof stipend, 5 points.}
The payment keeps the old basket affordable, so the maximum attainable
utility cannot fall. Reoptimisation at the new relative prices can raise
utility strictly above the old level. Hicksian compensation instead uses
$e(p^1,u^0)$, which is no greater than the new cost of the old basket.

Award 2 for affordability and the weak welfare comparison, 2 for the possible
strict gain, 1 for distinguishing the minimum expenditure benchmark.

\paragraph{2. The warranty and the raffle, 5 points.}
A strictly increasing globally concave wealth index rules out the independent ticket:
its expected net payoff is $-1/2$, so conditional Jensen makes it strictly
worse than not buying it. Loaded insurance can still be optimal, so the
contradiction comes from the ticket. Explaining both purchases requires
changing, for example, monetary-only payoffs, global concavity, or the
independence assumption.

Award 2 for the expected payoff and Jensen logic, 2 for identifying the
ticket rather than condemning insurance, 1 for a specific assumption change.
Simply calling both decisions irrational earns no credit. If T is the
ticket's net payoff, the formal argument is
\[
 \mathbb E[b(W+T)\mid W]\le b(W+\mathbb E T)=b(W-1/2)<b(W),
\]
using independence, positive final wealth and increasing concavity.
